% This is text associated with the open textbook "Open Electromagnetics"
% See immediately below for author, date
% This work is licensed under the Creative Commons Attribution-ShareAlike 4.0 International License. (CC BY SA 4.0) 
% To view a copy of the license, visit http://creativecommons.org/licenses/by-sa/4.0/

%ORB TITLE Wave Power in a Lossy Medium
%ORB AUTHOR 0        # S.W. Ellingson, Virginia Tech (ellingson@vt.edu)
%ORB DATE 20191115   # yyyymmdd
%ORB ABSTRACT  

\label{m0133_Wave_Power_in_a_Lossy_Medium} % <-- Must match name of this file and module directory name.  Don't mess with this!
\index{waves!power density|(}

In this section, we consider the power associated with waves propagating in materials which are potentially lossy; i.e., having conductivity $\sigma$ significantly greater than zero.
This topic has previously been considered in the section ``Wave Power in a Lossless Medium'' for the case in loss is not significant.\footnote{Depending on the version of this book, this section may appear in another volume.}
A review of that section may be useful before reading this section.

Recall that the Poynting vector
\index{Poynting vector}
%
\begin{equation}
{\bf S} \triangleq {\bf E} \times {\bf H} 
\end{equation}
%
indicates the power density (i.e., W/m$^2$) of a wave and the direction of power flow.
This is ``instantaneous'' power, applicable to waves regardless of the way they vary with time.
Often we are interested specifically in waves which vary sinusoidally, and which subsequently may be represented as phasors.  In this case, the \emph{time-average} Poynting vector is 
%
\begin{equation}
{\bf S}_{ave} \triangleq \frac{1}{2} \mbox{Re} \left\{ \widetilde{\bf E} \times \widetilde{\bf H}^* \right\}
\label{m0133_ePVP}
\end{equation}
%
Further, we have already used this expression to find that the time-average power density for a sinusoidally-varying uniform plane wave in a lossless medium is simply
%
\begin{equation}
S_{ave} = \frac{\left|E_0\right|^2}{2\eta} ~~~ \mbox{(lossless case)}
\label{m0133_ePaveL}
\end{equation}
%
where $\left|E_0\right|$ is the peak (as opposed to RMS) magnitude of the electric field intensity phasor, and $\eta$ is the wave impedance.

Let us now use Equation~\ref{m0133_ePVP} to determine the expression corresponding to Equation~\ref{m0133_ePaveL} in the case of possibly-lossy media.
We may express the electric and magnetic field intensities of a uniform plane wave as
%
\begin{equation}
\widetilde{\bf E} = \hat{\bf x}E_0 e^{-\alpha z} e^{-j\beta z}
\label{m0133_eE}
\end{equation}
%
and 
%
\begin{equation}
\widetilde{\bf H} = \hat{\bf y}\frac{E_0}{\eta_c} e^{-\alpha z} e^{-j\beta z}
\end{equation}
%
where $\alpha$ and $\beta$ are the attenuation constant and phase propagation constant, respectively, and 
$\eta_c$ is the complex-valued wave impedance.
As written, these expressions describe a wave which is $+\hat{\bf x}$-polarized and propagates in the $+\hat{\bf z}$ direction.  
We make these choices for convenience only -- as long as the medium is homogeneous and isotropic, we expect our findings to apply regardless of polarization and direction of propagation. 
Applying Equation~\ref{m0133_ePVP}:
%
\begin{flalign}
{\bf S}_{ave} &= \frac{1}{2} \mbox{Re} \left\{ \widetilde{\bf E} \times \widetilde{\bf H}^* \right\} \nonumber \\
&= \frac{1}{2} ~ \hat{\bf z} ~ \mbox{Re} \left\{ \frac{\left|E_0\right|^2}{\eta_c^*} e^{-2\alpha z} \right\} \nonumber \\
&= \hat{\bf z} \frac{\left|E_0\right|^2}{2} ~ \mbox{Re} \left\{ \frac{1}{\eta_c^*} \right\} e^{-2\alpha z} 
\label{m0133_eSave1}
\end{flalign}
%
Because $\eta_c$ is complex-valued when the material is lossy, we must proceed with caution.
First, let us write $\eta_c$ explicitly in terms of its magnitude $\left|\eta_c\right|$ and phase $\psi_{\eta}$:
%
\begin{equation}
\eta \triangleq \left|\eta\right| e^{j\psi_{\eta}}
\end{equation}
%
Then:
%
\begin{flalign}
\eta_c^* &= \left|\eta_c\right| e^{-j\psi_{\eta}} \\
\left(\eta_c^*\right)^{-1} &= \left|\eta_c\right|^{-1} e^{+j\psi_{\eta}} \\
\mbox{Re}\left\{\left(\eta_c^*\right)^{-1}\right\} &= \left|\eta_c\right|^{-1} \cos{\psi_{\eta}} 
\end{flalign}
%
Then Equation~\ref{m0133_eSave1} may be written:
%
\begin{equation}
\boxed{
{\bf S}_{ave} = \hat{\bf z} \frac{\left|E_0\right|^2}{2\left|\eta_c\right|} ~ e^{-2\alpha z} ~ \cos{\psi_{\eta}} 
}
\label{m0133_eSave} 
\end{equation}
%
%%%%%%%%%%%%%%%%%%%%%%%%%%%%%%%%%%%%%%%%%%%%%%%
%ORB HIGHLIGHT
\begin{mdframed}[backgroundcolor=yellow!20,nobreak=true] 
The time-average power density of the plane wave described by Equation~\ref{m0133_eE} in a possibly-lossy material is given by Equation~\ref{m0133_eSave}. 
\end{mdframed}
%%%%%%%%%%%%%%%%%%%%%%%%%%%%%%%%%%%%%%%%%%%%%%%

As a check, note that this expression gives the expected result for lossless media; i.e., $\alpha=0$, $\left|\eta_c\right|=\eta$, and $\psi_{\eta}=0$.
We now see that the effect of loss is that power density is now proportional to $\left(e^{-\alpha z}\right)^2$, so, as expected, power density is proportional to the square of either $\left|{\bf E}\right|$ or $\left|{\bf H}\right|$.
The result further indicates a one-time scaling of the power density by a factor of 
%
\begin{equation}
\frac{\left|\eta\right|}{\left|\eta_c\right|}\cos\psi_{\eta} < 1
\end{equation}
%
relative to a medium without loss.

%%%%%%%%%%%%%%%%%%%%%%%%%%%%%%%%%%%%%%%%%%%%%%%
%ORB HIGHLIGHT
\begin{mdframed}[backgroundcolor=yellow!20,nobreak=true] 
The reduction in power density due to non-zero conductivity is proportional to a distance-dependent factor $e^{-2\alpha z}$ and an additional factor that depends on the magnitude and phase of $\eta_c$. 
\end{mdframed}
%%%%%%%%%%%%%%%%%%%%%%%%%%%%%%%%%%%%%%%%%%%%%%%

\index{waves!power density|(}

